\subsubsection{The harmonic part of the isoparametric classification}
\label{sec:harmonic-isoparametric-scope}

For harmonic isoparametric polynomials, the dimension threshold in
\cref{ex:harmonic-isoparametric-radial} selects the following
parameters from the classical list.
The Cartan--M\"unzner equations \eqref{eq:CM-euclidean} and the
dimension relation \eqref{eq:iso-parameters} give
\[
 \Delta F=0\ \Longleftrightarrow\ m_1=m_2=m,
 \qquad N=gm+2.
\]
For \(g=4\), the possible equal multiplicities follow from the
homogeneous list of Takagi--Takahashi \cite{TakagiTakahashi72},
Stolz's multiplicity theorem \cite{Stolz99}, and the FKM construction
\cite[Theorem~4.1]{FKM81}. The sextic multiplicity restriction is due
to Abresch \cite{Abresch83}; see also
\cite[Chapter~3]{CecilRyan15}.
The first two columns give the multiplicity data; the last applies
the radial criterion:
\[
\begin{array}{c|c|c}
 g&\text{possible }N&\text{cases satisfying }N\ge4g\\ \hline
 2&2m+2,\ m\ge1&N\ge8\\
 3&5,8,14,26&14,26\\
 4&6,10&\text{none}\\
 6&8,14&\text{none}
\end{array}
\]
The FKM restriction follows directly from the Clifford relations.
A symmetric Clifford system
\(P_0,\ldots,P_a\) on \(\R^{2l}\) has FKM multiplicities
\((a,l-a-1)\), so equality requires \(l=2a+1\).
If \(a\ge2\), the operator \(P_1P_2\) preserves
\(E_+=\ker(P_0-I)\) and squares to \(-I\). Since
\(\dim E_+=l\), this forces \(l\) to be even, a contradiction.
Thus the only equal FKM pair is \((1,1)\), realized by
\(a=1,l=3\); the homogeneous list contributes the additional
quartic pair \((2,2)\). The quartic and sextic rows lie below the stability threshold in
\cref{prop:iso-stability}.
